\chapter{Background and Related Work}\label{ch:2}

\section{Text-to-SQL: Overview and Pipeline}\label{sec:2-1}

Given a natural language question Q and a database schema S = \{T, C, R\}, where T is a set of tables, C a set of columns, and R a set of foreign-key relations, the Text-to-SQL task seeks a function that maps (Q, S) to an executable SQL query Y that correctly answers Q against the database \citep{dong2023,safdarian2025}. Historically, Text-to-SQL research has passed through three broad stages. Early rule-based and grammar-based systems from the database community achieved reasonable accuracy only within narrow, single-domain settings, at the cost of extensive hand-engineering \citep{qin2022}. The introduction of large annotated cross-domain benchmarks, most influentially Spider \citep{yu2018} - shifted the field toward pre-trained language model (PLM) based encoder-decoder architectures, which learned to jointly encode the question and schema and decode a SQL query, resolving roughly 80\% of Spider's cases while still struggling on its hardest queries \citep{liu2024}. The most recent stage is dominated by large language models (LLMs), which perform Text-to-SQL largely through prompting rather than task-specific fine-tuning, pushing execution accuracy considerably higher while introducing new practical constraints around prompt length and inference cost \citep{liu2024,dong2023}.

Across all three stages, a common pipeline has emerged: a retrieval or schema-linking stage narrows the full database schema to a relevant subset; a generation stage produces a candidate SQL query conditioned on the question and the retrieved schema; and increasingly in the LLM era, a correction stage regenerates or repairs the query based on execution feedback or self-critique \citep{maamari2024}. This thesis is concerned exclusively with the first stage.

\section{Schema Linking: Definitions and Role}\label{sec:2-2}

Schema linking is the sub-task of identifying which tables and columns of a database schema are relevant to a given natural language question \citep{lei2020,wang2020}. \citet{safdarian2025} formalize it as a function that maps a question and a schema graph to a connected sub-schema sufficient to express the SQL query answering that question. Two motivations for this sub-task recur throughout the literature. The first is computational and economic: real-world schemas can contain hundreds of tables and thousands of columns, more than fits inside a model's context window, and every additional schema token included in a prompt has a direct cost in latency and API spend \citep{glass2025,wang2025b}. The second is accuracy: even when a schema does fit in context, including irrelevant tables and columns has been shown to actively degrade generation quality, since models can mistakenly select or join against irrelevant columns \citep{dong2023,maamari2024}.

\section{Existing Schema Linking Approaches}\label{sec:2-3}

Prior work on schema linking falls into four broad families, corresponding to the six methods evaluated in this thesis: lexical/rule-based matching, embedding or ranking-based neural matching, LLM-based reasoning (in forward, backward, and bidirectional variants), and graph-based structural matching.

\subsection{Lexical / Rule-Based Methods}\label{sec:2-3-1}

The earliest schema-linking components relied on surface-level string matching between question tokens and schema names, typically implemented as n-gram overlap or edit-distance heuristics. IRNet \citep{guo2019}, for instance, augments its encoder with hand-crafted type vectors derived from exactly this kind of string matching. These heuristics are attractive for being training-free and essentially instantaneous, but their limitations are well documented. \citet{lei2020} benchmark an automatic n-gram matching annotator of this kind against their own manually corrected gold standard and report only moderate agreement - an F1 of 0.775 for column references, 0.801 for tables - with most of the shortfall traced to paraphrases (a column named elevation referred to in the question as ``altitude'') and partial mentions that simple matching rules do not generalize to. \citet{taniguchi2021} reach a similar conclusion from their own independent annotation effort, motivating human annotation specifically because lexical heuristics were judged too noisy to serve as ground truth on their own. This body of evidence is precisely why the present thesis treats lexical matching as a baseline lower bound (Method A) rather than a candidate for the gold standard itself, and why its programmatically extracted Tier 2 gold standard is instead derived by parsing the gold SQL query rather than by string-matching the question.

\subsection{Embedding-Based Methods}\label{sec:2-3-2}

A second family replaces exact string overlap with a learned or pre-trained notion of semantic similarity between question and schema text. \citet{li2023} train a cross-encoder that jointly scores every table and column against the question and retains only the top-ranked four tables and five columns per table, reporting that this ranking step alone provides most of the benefit of their overall RESDSQL system. \citet{wang2025} take a retrieval-first approach at a larger scale, using multi-round semantic retrieval to narrow a pool of candidate databases and schema items before a separate grounding stage identifies the precise tables and columns needed, specifically to handle databases with thousands of fields spread across many separate schemas. \citet{glass2025} occupy a middle position: their extractive linker reuses the token-level hidden representations of a decoder-only LLM as a similarity signal, rather than training a dedicated embedding model, giving fine-grained control over the precision-recall trade-off at inference time without a generative decoding step. The embedding-based method evaluated in this thesis (Method B) follows the spirit of this family - scoring schema items by embedding similarity to the question rather than by exact lexical overlap - while deliberately using an off-the-shelf sentence embedding model rather than a task-specific fine-tuned cross-encoder, in order to isolate what a lightweight, training-free semantic method buys over lexical matching without confounding the comparison with a fine-tuning budget the other methods do not receive.

\subsection{LLM-Based Methods (Forward, Backward, Bidirectional)}\label{sec:2-3-3}

The most active area of recent schema-linking research uses an LLM directly as the linking mechanism, and the literature splits fairly cleanly along the direction of reasoning involved. \emph{Forward methods} reason from the question toward the schema: \citet{dong2023} C3 system prompts an LLM in a zero-shot setting to rank tables and columns by relevance to the question and retains the top candidates, reporting AUC scores of 0.9588 for table recall and 0.9833 for column recall on Spider's development set - evidence that a single well-prompted LLM call can rival specialized rankers without any fine-tuning. \emph{Backward methods} instead reason from an initial, possibly imperfect SQL query back to the schema: \citet{yang2024} SQL-to-Schema framework prompts an LLM to draft a full SQL query against the complete schema first, then parses that draft to recover the specific tables and columns it references, on the reasoning that a model reasoning about what a query would need to do is a more reliable schema-selection signal than a model reasoning about what a question is asking about in isolation. \emph{Bidirectional methods} combine both signals: \citet{cao2024} RSL-SQL runs forward and backward linking in parallel and unions the results, and their own ablation shows this union outperforms either direction alone for GPT-4o, though for a weaker DeepSeek-based model, backward linking alone slightly outperforms the bidirectional union, since the added recall from combining directions also introduces noise a weaker model handles less gracefully - evidence that the value of combining directions is itself model-dependent, not a fixed property of the technique.

A further LLM-based direction goes beyond a single forward or backward pass and treats schema exploration as an iterative or agentic process. \citet{wang2025b} AutoLink dynamically expands a linked schema subset over several LLM-guided steps rather than committing to a single-shot prediction, reporting strict recall above 97\% on the BIRD development set specifically because iterative expansion allows the model to recover from an initially incomplete guess. This body of work also speaks directly to two design choices this thesis makes for its own LLM-based methods. First, on sampling: \citet{katsogiannismeimarakis2026} show that increasing the number of sampled LLM outputs and taking their union systematically increases recall at the cost of precision, and \citet{dong2023} report a comparable pattern when varying the number of self-consistency samples for C3 - both consistent with this thesis's choice to use a small, fixed self-consistency budget (k = 3) for its forward LLM method rather than treating sample count as an unbounded lever for chasing recall. Second, on few-shot demonstrations: \citet{katsogiannismeimarakis2026} find that a single demonstration example meaningfully helps, but that additional demonstrations beyond one help inconsistently across models, cautioning against assuming that more in-context examples reliably improve LLM-based linking.

\subsection{Graph-Based Methods}\label{sec:2-3-4}

The final family exploits a structural property of relational schemas that the other three largely ignore: schemas are graphs, with tables as nodes and foreign-key relations as edges. Two rather different design points populate this family. \citet{wang2020} RAT-SQL treats the schema graph as an input to a neural relation-aware self-attention mechanism, learning soft, jointly-trained representations of question tokens and schema elements that are biased by, but not strictly limited to, the schema's actual foreign-key structure; this places RAT-SQL at the boundary between the graph-based and embedding-based families, since its schema linking is ultimately a byproduct of a trained encoder rather than a symbolic graph operation. \citet{safdarian2025} SchemaGraphSQL takes the opposite, training-free position: it uses a single LLM call only to identify coarse source and destination tables, then applies classical deterministic shortest-path algorithms over the schema's foreign-key graph to recover the full connected sub-schema joining them, reporting that this purely algorithmic approach outperforms specialized fine-tuned and prompt-heavy linkers on BIRD, with the path-finding step itself adding under 15 milliseconds of latency on top of a single lightweight LLM call. RAT-SQL and its more recent graph-neural descendants - including ShadowGNN \citep{chen2021}, LGESQL \citep{cao2021}, and SADGA \citep{cai2021} - represent a mature but comparatively heavyweight line of work requiring task-specific training; SchemaGraphSQL's classical path-finding represents a newer, lighter-weight alternative. The graph-based method evaluated in this thesis (Method G) follows this second, training-free tradition, since it is the one most compatible with this thesis's scope of comparing methods without a dedicated fine-tuning budget for any single family.

\section{Evaluation Practices for Schema Linking in Prior Work}\label{sec:2-4}

A recurring theme across the works reviewed above is the absence of a shared evaluation convention, a point \citet{glass2025} make explicitly: some schema-linking papers report area under the ROC curve, others report precision and recall, and the definitions of ``correct'' against which those metrics are computed are not the same from paper to paper. The papers surveyed in this chapter illustrate the point concretely rather than abstractly. \citet{li2023} and \citet{dong2023} both report AUC for table and column recall. \citet{lei2020} and \citet{taniguchi2021} each report precision, recall, and F1, but against two independently constructed gold standards that need not agree with one another on borderline cases. \citet{cao2024} introduce Strict Recall Rate, a stricter all-or-nothing variant of recall computed at the level of an entire column combination rather than per-column. \citet{yang2024} instead report table-recall@4, a metric defined relative to a fixed candidate-set size rather than an open-ended prediction. \citet{glass2025} argue for recall-weighted F-beta specifically because, in their own end-to-end experiments, the F6 score was well correlated with downstream SQL generation accuracy, while precision alone was not similarly predictive, motivating their recommendation that schema-linking evaluation should weight recall heavily rather than treating precision and recall as equally important.

This inconsistency is not merely a matter of academic bookkeeping. Because each paper's reported numbers are computed against a metric and a gold standard chosen by that paper's own authors, cross-paper comparison of the kind this thesis undertakes, asking whether an LLM-based method meaningfully outperforms a lexical or embedding-based one cannot be done reliably using published numbers alone; the methods have to be re-run and re-scored under one shared protocol. This is precisely the gap this thesis's five-metric evaluation framework (precision, recall, F1, SRR, and F6) is designed to close: rather than picking one of these conventions as authoritative, it computes all of them for every method, against both a human-annotated (Tier 1) and a programmatically derived (Tier 2) gold standard, so that a method's ranking can be checked for robustness across metric and gold-standard choice rather than being an artifact of whichever single metric happened to be reported.

\section{Gaps in Existing Literature}\label{sec:2-5}

Three gaps emerge from this review, each of which this thesis is designed to address directly.

\begin{itemize}
\item
  First, no work surveyed here evaluates lexical, embedding-based, LLM-based, and graph-based schema linking methods against one another under a single protocol; each paper instead introduces one new method and compares it to a small set of baselines drawn from whichever prior papers are most convenient, typically from only one or two of the four families.
\item
  Second, as documented in Section~\ref{sec:2-4}, the metrics and gold-standard definitions used to support those comparisons are themselves inconsistent across papers, which means that even a careful reader assembling numbers from multiple papers by hand cannot be confident the resulting comparison is fair.
\item
  Third, cost and practicality are rarely reported alongside accuracy: most of the LLM-based papers reviewed here report execution accuracy or schema-linking F-scores, but few report the API cost, latency, or engineering complexity of the method being proposed, making it difficult for a practitioner to weigh a small accuracy gain against a potentially large cost increase.
\end{itemize}

\begin{quote}
The ``does schema linking still help'' debate discussed in Section~\ref{sec:2-2} is itself a symptom of this last gap: \citet{maamari2024} and the large-scale schema-linking papers that motivate the opposite conclusion \citep{wang2025b,wang2025} are not actually in direct disagreement so much as they are answering the question under different, unstated assumptions about schema scale and model capability.
\end{quote}
